\section{Lista de acciones: 30 acciones aplicables en los próximos 12 meses}

Para convertir la discusión de esta entrevista en una ruta ejecutable, esta sección presenta una matriz de acciones organizada por trimestre. Cada acción se puede verificar de forma independiente, para evitar que «la estrategia sea correcta pero la ejecución quede en el vacío».

\begin{longtable}{p{0.08\textwidth}p{0.14\textwidth}p{0.18\textwidth}p{0.50\textwidth}}
\toprule
Número & Ventana de tiempo & Área responsable & Acción de ejecución y criterio de aceptación \\
\midrule
\endfirsthead
\toprule
Número & Ventana de tiempo & Área responsable & Acción de ejecución y criterio de aceptación \\
\midrule
\endhead
01 & T1 & Entrenamiento de modelos & Establecer una plantilla de registro experimental en tres etapas: pre, mid y post; el criterio de aceptación es que cualquier versión del modelo pueda rastrearse hasta el snapshot de los datos de entrenamiento, los hiperparámetros y el reporte de evaluación. \\
02 & T1 & Entrenamiento de modelos & Construir tres niveles de estrategia de inferencia, fast/think/pro, para las tareas centrales; el criterio de aceptación es que la misma tarea tenga una curva estable de desempeño-costo en distintos presupuestos. \\
03 & T1 & Ingeniería de datos & Establecer una clasificación por niveles del origen del corpus (licencia pública/licencia comercial/privado interno); el criterio de aceptación es que una muestra aleatoria de 200 elementos permita rastrear el origen. \\
04 & T1 & Ingeniería de datos & Introducir un pipeline de filtrado automático para datos sintéticos; el criterio de aceptación es que la tasa de aprobación, la tasa de falsos rechazos y la concordancia con la revisión manual puedan monitorearse de manera continua. \\
05 & T1 & Ingeniería de plataforma & Establecer un experimento de programación con separación actor/learner; el criterio de aceptación es que la relación entre el aumento de rendimiento y la sobrecarga de comunicación pueda cuantificarse. \\
06 & T1 & Ingeniería de plataforma & Unificar el mecanismo de reversión al publicar cambios; el criterio de aceptación es que los cambios de alto riesgo puedan revertirse a una versión estable en 15 minutos. \\
07 & T1 & Ingeniería de producto & Introducir una plantilla de especificación para el coding agent; el criterio de aceptación es que disminuyan las rondas de aclaración de requisitos y aumente la tasa de aprobación en el primer intento. \\
08 & T1 & Ingeniería de producto & Introducir aceptación automatizada (pruebas unitarias, análisis estático, normas de estilo); el criterio de aceptación es que la tasa de aprobación del código enviado por la IA aumente y se mantenga estable. \\
09 & T1 & Gobernanza y asuntos legales & Establecer un mecanismo de clasificación de incidentes por salidas del modelo; el criterio de aceptación es que cada incidente tenga responsable asignado, tiempo de corrección y conclusiones de la revisión posterior. \\
10 & T1 & Gobernanza y asuntos legales & Introducir un proceso de respaldo humano para escenarios de alto riesgo; el criterio de aceptación es que todas las tareas sensibles cuenten con un registro de «confirmación final humana». \\
11 & T2 & Entrenamiento de modelos & Impulsar el entrenamiento especializado con RLVR en tareas verificables; el criterio de aceptación es que la precisión en la tarea objetivo mejore de manera significativa sin que se deteriore la generalización. \\
12 & T2 & Entrenamiento de modelos & Introducir conjuntos de evaluación multitarea divididos por rangos de dificultad; el criterio de aceptación es que el progreso del modelo no dependa de un único benchmark. \\
13 & T2 & Ingeniería de datos & Construir un monitoreo de contaminación de datos (proporción generada por IA, tasa de duplicados, deriva de calidad); el criterio de aceptación es que las fluctuaciones anómalas generen alertas dentro de 24 horas. \\
14 & T2 & Ingeniería de datos & Establecer un punto de referencia de calidad para OCR y extracción de documentos; el criterio de aceptación es que la precisión y la exhaustividad de la extracción estructurada puedan rastrearse con ambos indicadores. \\
15 & T2 & Ingeniería de plataforma & Construir una cadena de recuperación-compresión-retroceso para contexto largo; el criterio de aceptación es que disminuya la proporción de fallos en tareas complejas causados por «pérdida de información». \\
16 & T2 & Ingeniería de plataforma & Impulsar la conversión en producto del servicio de enrutamiento de modelos; el criterio de aceptación es que la estrategia de cambio de modelo sea configurable, admita despliegue gradual y pueda reproducirse. \\
17 & T2 & Ingeniería de producto & Aplicar el principio de mínimo privilegio a los permisos de las herramientas; el criterio de aceptación es que cada llamada a una herramienta tenga un registro de auditoría y una anotación de uso. \\
18 & T2 & Ingeniería de producto & Establecer un ciclo cerrado que lleve la retroalimentación de los usuarios a muestras de entrenamiento; el criterio de aceptación es que el tiempo de procesamiento de la retroalimentación y la tasa de adopción mejoren de manera sostenida. \\
19 & T2 & Gobernanza y asuntos legales & Completar la auditoría anual de licencias de datos; el criterio de aceptación es que el corpus de alto riesgo se elimine o se reemplace por completo, y que el reporte de evaluación de riesgos pueda revisarse. \\
20 & T2 & Gobernanza y asuntos legales & Establecer una revisión de seguridad previa al lanzamiento del modelo; el criterio de aceptación es que cada lanzamiento cuente con un registro de revisión completo. \\
21 & T3 & Entrenamiento de modelos & Probar esquemas de puente semiverificable para «tareas no verificables»; el criterio de aceptación es que al menos un tipo de tarea forme una ruta de mejora estable. \\
22 & T3 & Entrenamiento de modelos & Añadir entrenamiento impulsado por muestras de fallo (hard case replay); el criterio de aceptación es que disminuya la tasa de recurrencia de problemas de fallo frecuentes históricos. \\
23 & T3 & Ingeniería de datos & Establecer un mecanismo de actualización incremental de la base de conocimiento privada; el criterio de aceptación es que el ciclo para que el nuevo conocimiento entre al sistema se acorte de manera notable. \\
24 & T3 & Ingeniería de plataforma & Evaluar la ruta de hardware con separación de entrenamiento e inferencia; el criterio de aceptación es que se forme un plan ejecutable de costo de hardware y migración de desempeño. \\
25 & T3 & Ingeniería de producto & Establecer un protocolo de colaboración multiagente (división de tareas, manejo de conflictos); el criterio de aceptación es que disminuya la tasa de fallos en la colaboración entre tareas. \\
26 & T3 & Gobernanza y asuntos legales & Introducir un mecanismo de escalamiento para conversaciones de alto riesgo (intervención humana); el criterio de aceptación es que la tasa de manejo incorrecto en escenarios sensibles disminuya de manera sostenida. \\
27 & T4 & Entrenamiento de modelos & Realizar una revisión posterior de la distribución del presupuesto anual de entrenamiento; el criterio de aceptación es que se establezca una relación explicable entre la inversión de presupuesto y el beneficio de negocio. \\
28 & T4 & Ingeniería de plataforma & Completar el sistema de observabilidad de todo el año; el criterio de aceptación es que disminuyan tanto el tiempo de localización de fallas en producción como el tiempo de recuperación. \\
29 & T4 & Ingeniería de producto & Completar la capacitación interna en «habilidades de colaboración humano-máquina»; el criterio de aceptación es que el equipo muestre una mejora cuantificable en la redacción de especificaciones y la capacidad de revisión. \\
30 & T4 & Gobernanza y asuntos legales & Publicar el reporte anual de responsabilidad de IA; el criterio de aceptación es que se divulgue públicamente el marco de gobernanza, las estadísticas de incidentes y el cierre del ciclo de acciones de mejora. \\
\bottomrule
\end{longtable}

\begin{importantbox}{Recomendación de prioridad de ejecución}
Si los recursos son limitados, conviene priorizar tres cosas: primero, el ciclo cerrado de evaluación verificable; segundo, el ciclo cerrado de auditoría del origen de los datos; tercero, el ciclo cerrado de reversión al publicar cambios. Mientras estos tres ciclos cerrados se sostengan, la ruta del modelo puede iterar rápidamente por prueba y error sin que el sistema pierda el control.
\end{importantbox}

\begin{knowledgebox}{Consenso mínimo a nivel organizacional}
Cualquier «estrategia de IA» requiere la firma conjunta de tres partes: el responsable técnico (viabilidad), el responsable de producto (valor) y el responsable legal/de cumplimiento (sostenibilidad). Si falta cualquiera de las partes, la implementación enfrentará resistencia estructural más adelante.
\end{knowledgebox}

\subsection{Resumen de la sección}

La discusión de la ruta debe traducirse en unidades ejecutables. Descomponer los juicios generales en acciones verificables es la única forma confiable de evitar que la estrategia quede en el vacío.

